\chapter{2009年湖北卷（文）}

\section{单选题}

\begin{problem}
若向量 \(\bs{a} = \left(1,1\right)\)，\(\bs{b} = \left(-1,1\right)\)，\(\bs{c} = \left(4,2\right)\)，则 \(\bs{c} =\)（\quad）

\choices
    {\(3\bs{a} + \bs{b}\)}
    {\(3\bs{a} - \bs{b}\)}
    {\(-\bs{a} + 3\bs{b}\)}
    {\(\bs{a} + 3\bs{b}\)}
\end{problem}

\begin{answer}
B
\end{answer}

\begin{solution}
计算得 \(3\bs{a}-\bs{b}=3(1,1)-(-1,1)=(3,3)-(-1,1)=(4,2)=\bs{c}\)，所以选 B.
\end{solution}

\begin{problem}
函数 \( y = \frac{1 - 2x}{1 + 2x} \)（\(x \in \R\)，\(x \neq -\frac{1}{2}\)）的反函数是（\quad）

\choices
    {\( y = \frac{1 + 2x}{1 - 2x} \)（\(x \in \R\)，\(x \neq \frac{1}{2}\)）}
    {\( y = \frac{1 - 2x}{1 + 2x} \)（\(x \in \R\)，\(x \neq -\frac{1}{2}\)）}
    {\( y = \frac{1 + x}{2\left(1 - x\right)} \)（\(x \in \R\)，\(x \neq 1\)）}
    {\( y = \frac{1 - x}{2\left(1 + x\right)} \)（\(x \in \R\)，\(x \neq -1\)）}
\end{problem}

\begin{answer}
D
\end{answer}

\begin{solution}
由 \(y=\dfrac{1-2x}{1+2x}\) 得 \(y(1+2x)=1-2x\)，所以 \(2x(y+1)=1-y\)，从而 \(x=\dfrac{1-y}{2(1+y)}\)，且 \(y\neq-1\). 交换 \(x\)，\(y\) 得反函数为 \(y=\dfrac{1-x}{2(1+x)}\)，定义域为 \(x\in\R\) 且 \(x\neq-1\)，故选 D.
\end{solution}

\begin{problem}
“\(\sin \alpha = \frac{1}{2}\)”是“\(\cos 2\alpha = \frac{1}{2}\)”的（\quad）

\choices
    {充分而不必要条件}
    {必要而不充分条件}
    {充要条件}
    {既不充分也不必要条件}
\end{problem}

\begin{answer}
A
\end{answer}

\begin{solution}
若 \(\sin\alpha=\dfrac{1}{2}\)，则由二倍角公式 \(\cos2\alpha=1-2\sin^2\alpha=1-2\times\dfrac{1}{4}=\dfrac{1}{2}\)，因此前者是后者的充分条件. 反过来，若 \(\cos2\alpha=\dfrac{1}{2}\)，则 \(\sin^2\alpha=\dfrac{1}{4}\)，只能得到 \(\sin\alpha=\pm\dfrac{1}{2}\)，例如 \(\alpha=-\dfrac{\pi}{6}\) 时不满足 \(\sin\alpha=\dfrac{1}{2}\)，所以不是必要条件，故选 A.
\end{solution}

\begin{problem}
从 \(5\) 名志愿者中选派 \(4\) 人在星期五、星期六、星期日参加公益活动，每人一天，要求星期五有一人参加，星期六有两人参加，星期日有一人参加，则不同的选派方法共有（\quad）

\choices
    {\(120\text{ 种}\)}
    {\(96\text{ 种}\)}
    {\(60\text{ 种}\)}
    {\(48\text{ 种}\)}
\end{problem}

\begin{answer}
C
\end{answer}

\begin{solution}
先从 \(5\) 人中选出星期五参加的 \(1\) 人，有 \(5\) 种选法；再从余下 \(4\) 人中选出星期六参加的 \(2\) 人，有 \(\mathrm C_4^2=6\) 种选法；最后从剩下的 \(2\) 人中选出星期日参加的 \(1\) 人，有 \(2\) 种选法. 因而不同的选派方法共有 \(5\times6\times2=60\) 种，故选 C.
\end{solution}

\begin{problem}
已知双曲线 \(\frac{x^2}{2} - \frac{y^2}{2} = 1\) 的准线经过椭圆 \(\frac{x^2}{4} + \frac{y^2}{b^2} = 1\)（\(b > 0\)）的焦点，则（\quad）

\choices
    {\(3\)}
    {\(\sqrt{5}\)}
    {\(\sqrt{3}\)}
    {\(\sqrt{2}\)}
\end{problem}

\begin{answer}
C
\end{answer}

\begin{solution}
双曲线可写成 \(\dfrac{x^2}{a^2}-\dfrac{y^2}{d^2}=1\)，其中 \(a^2=d^2=2\)，焦半距 \(c_1=\sqrt{a^2+d^2}=2\)，离心率 \(e=\dfrac{c_1}{a}=\sqrt2\)，所以其准线为 \(x=\pm\dfrac{a}{e}=\pm1\). 椭圆的焦点在 \(x\) 轴上，焦点横坐标的绝对值为 \(\sqrt{4-b^2}\). 准线经过焦点，故 \(\sqrt{4-b^2}=1\)，于是 \(b^2=3\). 由 \(b>0\) 得 \(b=\sqrt3\)，故选 C.
\end{solution}

\begin{problem}
如图，在三棱柱 \(ABC - A_{1}B_{1}C_{1}\) 中，\(\angle ACB = 90^{\circ}\)，\(\angle ACC_{1} = 60^{\circ}\)，\(\angle BCC_{1} = 45^{\circ}\)，侧棱 \(CC_{1}\) 的长为 \(1\)，则该三棱柱的高等于（\quad）

\bitmapfigure[max width=0.34\linewidth]{img/2009/hubei_liberal/q6_fig1.png}

\choices
    {\(\frac{1}{2}\)}
    {\(\frac{\sqrt{2}}{2}\)}
    {\(\frac{\sqrt{3}}{2}\)}
    {\(\frac{\sqrt{3}}{3}\)}
\end{problem}

\begin{answer}
A
\end{answer}

\begin{solution}
取底面内互相垂直的单位方向分别沿 \(CA\)，\(CB\)，再取垂直于底面的单位方向. 设 \(CC_1\) 在这三个方向上的分量分别为 \(u\)，\(v\)，\(h\)，其中 \(h\) 是三棱柱的高. 因为 \(CC_1=1\)，且给出的两个夹角分别为 \(60^\circ\)，\(45^\circ\)，所以 \(u=\cos60^\circ=\dfrac12\)，\(v=\cos45^\circ=\dfrac{\sqrt2}{2}\). 由勾股关系，
\[
h^2=1-u^2-v^2=1-\frac14-\frac12=\frac14
\]
由于 \(h>0\)，得 \(h=\dfrac12\)，故选 A.
\end{solution}

\begin{problem}
函数 \( y = \cos \left(2x + \frac{\pi}{6}\right) - 2 \) 的图象 \(F\) 按向量 \(\bs{a}\) 平移到 \(F'\)，\(F'\) 的解析式 \( y = f(x) \)，当 \( y = f(x) \) 为奇函数时，向量 \(\bs{a}\) 可以等于（\quad）

\choices
    {\(\left(\frac{\pi}{6}, -2\right)\)}
    {\(\left(\frac{\pi}{6}, 2\right)\)}
    {\(\left(-\frac{\pi}{6}, -2\right)\)}
    {\(\left(-\frac{\pi}{6}, 2\right)\)}
\end{problem}

\begin{answer}
D
\end{answer}

\begin{solution}
设平移向量为 \(\bs{a}=(h,k)\). 平移后图象的函数为 \(f(x)=\cos\left(2(x-h)+\dfrac{\pi}{6}\right)-2+k\). 记 \(\phi=-2h+\dfrac{\pi}{6}\)，则 \(f(x)=\cos(2x+\phi)+k-2\). 若 \(f\) 为奇函数，则 \(f(x)+f(-x)\) 恒为零，而
\[
f(x)+f(-x)=2\cos2x\cos\phi+2(k-2)
\]
因此必须有 \(\cos\phi=0\) 且 \(k=2\). 此时
\[
f(x)=\cos\left(2x-2h+\frac{\pi}{6}\right)
\]
要使这个余弦函数为奇函数，必须有 \(\cos\left(-2h+\dfrac{\pi}{6}\right)=0\)，即 \(-2h+\dfrac{\pi}{6}=\dfrac{\pi}{2}+m\pi\)，其中 \(m\) 为整数. 在给出的选项中，只有 \(h=-\dfrac{\pi}{6}\) 满足这一条件；此时 \(f(x)=\cos\left(2x+\dfrac{\pi}{2}\right)=-\sin2x\)，确为奇函数. 因而 \(\bs{a}=\left(-\dfrac{\pi}{6},2\right)\)，故选 D.
\end{solution}

\begin{problem}
在“家电下乡”活动中，某厂要将 \(100\text{ 台}\)洗衣机运往邻近的乡镇，现有 \(4\text{ 辆}\)甲型货车和 \(8\text{ 辆}\)乙型货车可供使用，每辆甲型货车运输费用 \(400\text{ 元}\)，可装洗衣机 \(20\text{ 台}\)；每辆乙型货车运输费用 \(300\text{ 元}\)，可装洗衣机 \(10\text{ 台}\)，若每辆至多只运一次，则该厂所花的最少运输费用为（\quad）

\choices
    {\(2000\text{ 元}\)}
    {\(2200\text{ 元}\)}
    {\(2400\text{ 元}\)}
    {\(2800\text{ 元}\)}
\end{problem}

\begin{answer}
B
\end{answer}

\begin{solution}
设使用甲型货车 \(x\) 辆、乙型货车 \(y\) 辆，则 \(x\)，\(y\) 为非负整数，且 \(0\leq x\leq4\)，\(0\leq y\leq8\)，\(20x+10y\geq100\)，即 \(2x+y\geq10\). 因为 \(y\leq8\)，必须有 \(x\geq1\). 对固定的 \(x=1\)，\(2\)，\(3\)，\(4\)，费用最少时分别取 \(y=8\)，\(6\)，\(4\)，\(2\)，对应费用为 \(2800\)，\(2600\)，\(2400\)，\(2200\) 元. 因此最少运输费用为 \(2200\) 元，故选 B.
\end{solution}

\begin{problem}
设 \(x \in \R\)，记不超过 \(x\) 的最大整数为 \([x]\)，令 \(\{x\} = x - [x]\)，则 \(\left\{\frac{\sqrt{5}+1}{2}\right\}\)，\(\left[\frac{\sqrt{5}+1}{2}\right]\)，\(\frac{\sqrt{5}+1}{2}\)（\quad）

\choices
    {是等差数列但不是等比数列}
    {是等比数列但不是等差数列}
    {既是等差数列又是等比数列}
    {既不是等差数列也不是等比数列}
\end{problem}

\begin{answer}
B
\end{answer}

\begin{solution}
令 \(t=\dfrac{\sqrt5+1}{2}\)，则 \(1<t<2\)，所以 \([t]=1\)，\(\{t\}=t-1=\dfrac{\sqrt5-1}{2}\). 又因为 \(t^2=t+1\)，故 \(t(t-1)=1\)，从而三项 \(t-1\)，\(1\)，\(t\) 的相邻两项之比均为 \(t\)，它们构成等比数列. 若它们构成等差数列，则应有 \(1-(t-1)=t-1\)，即 \(t=\dfrac32\)，但 \(t=\dfrac{\sqrt5+1}{2}\neq\dfrac32\)，所以不是等差数列，故选 B.
\end{solution}

\begin{problem}
古希腊人常用小石子在沙滩上摆成各种形状来研究数，例如：他们研究过图1中的\(1\)，\(3\)，\(6\)，\(10\)，\(\cdots\)，由于这些数能够表示成三角形，将其称为三角形数；类似地，称图2中的\(1\)，\(4\)，\(9\)，\(16\)，\(\cdots\)这样的数称为正方形数. 下列数中既是三角形数又是正方形数的是（\quad）

\examfiguregroup[float=inline]{img/2009/hubei_liberal/q10_fig1.png}{%
    \vbox{%
        \hbox to \linewidth{%
            \hfil\bitmapinclude[max width=0.28\linewidth]{img/2009/hubei_liberal/q10_fig1.png}\hfil
        }%
        \vskip 0.4\baselineskip
        \hbox to \linewidth{%
            \hfil\bitmapinclude[max width=0.30\linewidth]{img/2009/hubei_liberal/q10_fig2.png}\hfil
        }%
    }%
}
\FloatBarrier

\choices
    {\(289\)}
    {\(1024\)}
    {\(1225\)}
    {\(1378\)}
\end{problem}

\begin{answer}
C
\end{answer}

\begin{solution}
第 \(n\) 个三角形数为 \(T_n=\dfrac{n(n+1)}{2}\)，正方形数必须是完全平方数. \(289=17^2\)，但若它是三角形数，则应有 \(n(n+1)=578\)，而 \(23\times24=552<578<600=24\times25\)，故 \(289\) 不是三角形数. \(1024=32^2\)，而 \(T_{44}=\dfrac{44\times45}{2}=990\)，\(T_{45}=\dfrac{45\times46}{2}=1035\)，故不是三角形数. \(1225=35^2\)，且 \(T_{49}=\dfrac{49\times50}{2}=1225\)，所以它同时是两类数. 最后 \(37^2=1369<1378<1444=38^2\)，故 \(1378\) 不是正方形数. 因此应选 C.
\end{solution}

\section{填空题}

\begin{problem}
已知 \( (1 + ax)^{5} = 1 + 10x + bx^{2} + \cdots + a^{5}x^{5} \)，则 \(b=\fillinblank{}\).
\end{problem}

\begin{answer}
40
\end{answer}

\begin{solution}
由二项式定理，\((1+ax)^5\) 中 \(x\) 的系数为 \(5a\)，所以 \(5a=10\)，得 \(a=2\). \(x^2\) 的系数为 \(\mathrm C_5^2a^2=10\times2^2=40\)，故 \(b=40\).
\end{solution}

\begin{problem}
甲、乙、丙三人将参加某项测试，他们能达标的概率分别是 \(0.8\)，\(0.6\)，\(0.5\)，则三人都达标的概率是\fillinblank{}，三人中至少有一人达标的概率是\fillinblank{}.
\end{problem}

\begin{answer}
0.24，0.96
\end{answer}

\begin{solution}
题面只给出了三人的边际概率，按此类题的通常约定将三人的达标事件视为相互独立；若不作独立性约定，第二空不能由题面唯一确定. 在该约定下，三人都达标的概率为 \(0.8\times0.6\times0.5=0.24\). 三人都不达标的概率为 \((1-0.8)(1-0.6)(1-0.5)=0.2\times0.4\times0.5=0.04\)，所以至少有一人达标的概率为 \(1-0.04=0.96\).
\end{solution}

\begin{problem}
设集合 \(A = \{x \mid \log_2 x < 1\}\)，\(B = \left\{x \mid \frac{x - 1}{x + 2} < 1\right\}\)，则 \(A \cap B=\fillinblank{}\).
\end{problem}

\begin{answer}
\(\{x \mid 0 < x < 2\}\)
\end{answer}

\begin{solution}
由对数式有 \(x>0\)，且 \(\log_2x<1\) 等价于 \(x<2\)，所以 \(A=\{x\mid0<x<2\}\). 对 \(B\)，先注意 \(x\neq-2\)，并移项得 \(\dfrac{-3}{x+2}<0\)，故 \(x+2>0\)，即 \(B=\{x\mid x>-2\}\). 因而 \(A\cap B=\{x\mid0<x<2\}\).
\end{solution}

\begin{problem}
过原点 \(O\) 作圆 \( x^{2} + y^{2} - 6x - 8y + 20 = 0 \) 的两条切线，设切点分别为 \(P\)，\(Q\)，则线段 \(PQ\) 的长为\fillinblank{}.
\end{problem}

\begin{answer}
4
\end{answer}

\begin{solution}
将圆的方程配方得 \((x-3)^2+(y-4)^2=5\)，所以圆心为 \(C(3,4)\)，半径为 \(\sqrt5\)，且 \(OC=5\). 设 \(P(x,y)\) 为切点，则 \(CP\perp OP\)，即 \((P-C)\cdot P=0\)，从而 \(x^2+y^2-3x-4y=0\). 与圆的方程相减得切点弦 \(PQ\) 的方程 \(3x+4y=20\). 圆心到该弦的距离为 \(\dfrac{|3\times3+4\times4-20|}{\sqrt{3^2+4^2}}=1\)，所以
\[
PQ=2\sqrt{(\sqrt5)^2-1^2}=4
\]
\end{solution}

\begin{problem}
下图是样本容量为 \(200\) 的频率分布直方图.

\texfigure[max width=0.36\linewidth]{img_repaint/2009/hubei_liberal/q15_fig1.tex}

根据样本的频率分布直方图估计，样本数据落在 \([6,10)\) 内的频数为\fillinblank{}，数据落在 \([2,10)\) 内的概率约为\fillinblank{}.
\end{problem}

\begin{answer}
64，0.4
\end{answer}

\begin{solution}
频率分布直方图中，某组的频率等于组距乘以该组的频率组距. 区间 \([6,10)\) 的组距为 \(4\)，高度为 \(0.08\)，所以其频率为 \(4\times0.08=0.32\)，频数为 \(200\times0.32=64\). 区间 \([2,10)\) 包含 \([2,6)\) 和 \([6,10)\)，其概率约为 \(4\times0.02+4\times0.08=0.40\).
\end{solution}

\section{解答题}

\begin{problem}
在锐角 \(\triangle ABC\) 中，\(a\)，\(b\)，\(c\) 分别为角 \(A\)，\(B\)，\(C\) 所对的边，且 \(\sqrt{3} a = 2c \sin A\).

\begin{enumerate}
\item 确定角 \(C\) 的大小；
\item 若 \(c = \sqrt{7}\)，且 \(\triangle ABC\) 的面积为 \(\frac{3\sqrt{3}}{2}\)，求 \(a + b\) 的值.
\end{enumerate}
\end{problem}

\begin{solution}
\begin{enumerate}
\item 由正弦定理，\(\dfrac{a}{\sin A}=\dfrac{c}{\sin C}\)，所以 \(\sin A=\dfrac{a\sin C}{c}\). 代入已知条件 \(\sqrt3a=2c\sin A\)，得 \(\sqrt3a=2a\sin C\). 由于 \(a>0\)，所以 \(\sin C=\dfrac{\sqrt3}{2}\). 又三角形为锐角三角形，\(0<C<\dfrac{\pi}{2}\)，故 \(C=\dfrac{\pi}{3}\).
\item 由三角形面积公式，\(\dfrac12ab\sin C=\dfrac{3\sqrt3}{2}\)，从而 \(\dfrac12ab\times\dfrac{\sqrt3}{2}=\dfrac{3\sqrt3}{2}\)，得 \(ab=6\). 再由余弦定理，\(c^2=a^2+b^2-2ab\cos C\)，所以 \(7=a^2+b^2-2\times6\times\dfrac12=a^2+b^2-6\)，即 \(a^2+b^2=13\). 因而 \((a+b)^2=a^2+b^2+2ab=13+12=25\). 由于 \(a+b>0\)，得 \(a+b=5\).
\end{enumerate}
\end{solution}

\begin{problem}
围建一个面积为 \(360 \mathrm{~m}^{2}\) 的矩形场地，要求矩形场地的一面利用旧墙（利用旧墙需维修），其它三面围墙要新建. 在旧墙的对面的新墙上要留一个宽度为 \(2 \mathrm{~m}\) 的进出口. 如图所示，已知旧墙的维修费用为 \(45\text{ 元}/\mathrm{m}\)，新墙的造价为 \(180\text{ 元}/\mathrm{m}\). 设利用的旧墙的长度为 \(x\)（单位：\(\mathrm{m}\)），修建此矩形场地围墙的总费用为 \(y\)（单位：\text{元}）.

\begin{enumerate}
\item 将 \(y\) 表示为 \(x\) 的函数；
\item 试确定 \(x\)，使修建此矩形场地围墙的总费用最少，并求出最小总费用.
\end{enumerate}

\FigureLayoutDeclare{img/2009/hubei_liberal/q17_fig1.png}{10.0cm}{29bp 44bp 26bp 31bp}
\bitmapfigure[max width=0.42\linewidth]{img/2009/hubei_liberal/q17_fig1.png}
\end{problem}

\begin{solution}
\begin{enumerate}
\item 设与旧墙垂直的两面新墙长度均为 \(z\)，则 \(xz=360\)，所以 \(z=\dfrac{360}{x}\). 旧墙维修费为 \(45x\) 元；旧墙对面的新墙扣除 \(2\mathrm{~m}\) 进出口后，造价为 \(180(x-2)\) 元；另外两面新墙造价为 \(2\times180z\) 元. 因而 \(x>2\)，且总费用为
\[
y=45x+180(x-2)+360z=225x+\frac{129600}{x}-360
\]
\item 由基本不等式，
\[
225x+\frac{129600}{x}\geq2\sqrt{225x\cdot\frac{129600}{x}}=10800
\]
当且仅当 \(225x=\dfrac{129600}{x}\)，即 \(x^2=576\)，又 \(x>2\)，所以 \(x=24\). 此时总费用的最小值为 \(10800-360=10440\) 元.
\end{enumerate}
\end{solution}

\begin{problem}
如图，四棱锥 \(S - ABCD\) 的底面是正方形，\(SD \perp\) 平面 \(ABCD\)，\(SD = AD = a\)，点 \(E\) 是 \(SD\) 上的点，且 \(DE = \lambda a\)，\(\lambda \in \left(0,1\right]\).

\begin{enumerate}
\item 求证：对任意的 \(\lambda \in \left(0,1\right]\)，都有 \(AC \perp BE\)；
\item 若二面角 \(C - AE - D\) 的大小为 \(60^{\circ}\)，求 \(\lambda\) 的值.
\end{enumerate}

\bitmapfigure[max width=0.34\linewidth]{img/2009/hubei_liberal/q18_fig1.png}
\end{problem}

\begin{solution}
\begin{enumerate}
\item 建立空间直角坐标系，取 \(D\) 为原点，令 \(DA\)，\(DC\)，\(DS\) 分别为 \(x\)，\(y\)，\(z\) 轴正方向. 由题意可设 \(D(0,0,0)\)，\(A(a,0,0)\)，\(C(0,a,0)\)，\(B(a,a,0)\)，\(S(0,0,a)\). 因为 \(DE=\lambda a\)，所以 \(E(0,0,\lambda a)\). 于是 \(\overrightarrow{AC}=(-a,a,0)\)，\(\overrightarrow{BE}=(-a,-a,\lambda a)\)，从而
\[
\overrightarrow{AC}\cdot\overrightarrow{BE}=a^2-a^2+0=0
\]
因此 \(AC\perp BE\).
\item 平面 \(CAE\) 与平面 \(DAE\) 的交线为 \(AE\). 取这两个平面的法向量
\(\bs{n}_1=\overrightarrow{AC}\times\overrightarrow{AE}=a^2(\lambda,\lambda,1)\)，\(\bs{n}_2=(0,1,0)\)
其中 \(\bs{n}_2\) 是平面 \(DAE\) 的一个法向量. 因为 \(0<\lambda\leq1\)，二面角的大小为锐角，故其余弦为
\[
\cos60^\circ=\frac{\bs{n}_1\cdot\bs{n}_2}{|\bs{n}_1||\bs{n}_2|}=\frac{\lambda}{\sqrt{1+2\lambda^2}}
\]
于是 \(\dfrac12=\dfrac{\lambda}{\sqrt{1+2\lambda^2}}\)，平方并整理得 \(4\lambda^2=1+2\lambda^2\)，即 \(\lambda^2=\dfrac12\). 结合 \(\lambda>0\)，得 \(\lambda=\dfrac{\sqrt2}{2}\).
\end{enumerate}
\end{solution}

\begin{problem}
已知 \(\{a_{n}\}\) 是一个公差大于 \(0\) 的等差数列，且满足 \(a_{3}a_{6} = 55\)，\(a_{2} + a_{7} = 16\).

\begin{enumerate}
\item 求数列 \(\{a_{n}\}\) 的通项公式；
\item 若数列 \(\{a_{n}\}\) 和数列 \(\{b_{n}\}\) 满足等式：\(a_{n} = \frac{b_{1}}{2} + \frac{b_{2}}{2^{2}} + \frac{b_{3}}{2^{3}} + \cdots + \frac{b_{n}}{2^{n}}\)（\(n\) 为正整数）. 求数列 \(\{b_{n}\}\) 的前 \(n\) 项和 \(S_{n}\).
\end{enumerate}
\end{problem}

\begin{solution}
\begin{enumerate}
\item 设等差数列的首项为 \(a_1\)，公差为 \(d\)，则 \(a_3=a_1+2d\)，\(a_6=a_1+5d\). 又 \(a_2+a_7=(a_1+d)+(a_1+6d)=a_3+a_6=16\). 因此 \(a_3\)，\(a_6\) 是方程 \(t^2-16t+55=0\) 的两个根，即 \((t-5)(t-11)=0\). 因为 \(d>0\)，所以 \(a_6>a_3\)，从而 \(a_3=5\)，\(a_6=11\). 于是 \(3d=a_6-a_3=6\)，得 \(d=2\)，再由 \(a_3=a_1+2d\) 得 \(a_1=1\). 所以数列通项为 \(a_n=a_1+(n-1)d=2n-1\).
\item 当 \(n=1\) 时，题给关系为 \(a_1=\dfrac{b_1}{2}\)，所以 \(b_1=2\). 当 \(n\geq2\) 时，用第 \(n\) 项关系减去第 \(n-1\) 项关系，得 \(a_n-a_{n-1}=\dfrac{b_n}{2^n}\). 由于 \(a_n-a_{n-1}=2\)，所以 \(b_n=2^{n+1}\). 因此
\[
S_n=b_1+\sum_{k=2}^n b_k=2+\sum_{k=2}^n2^{k+1}=2+\left(2^3+2^4+\cdots+2^{n+1}\right)=2^{n+2}-6
\]
当 \(n=1\) 时右式也等于 \(2\)，故公式对所有正整数 \(n\) 成立.
\end{enumerate}
\end{solution}

\begin{problem}
如图，过抛物线 \( y^{2} = 2px \)（\(p > 0\)）的焦点 \(F\) 的直线与抛物线相交于 \(M\)，\(N\) 两点，自 \(M\)，\(N\) 向准线 \( l \) 作垂线，垂足分别为 \(M_{1}\)，\(N_{1}\).

\begin{enumerate}
\item 求证：\( FM_{1} \perp FN_{1} \)；
\item 记 \(\triangle FMM_{1}\)，\(\triangle FM_{1}N_{1}\)，\(\triangle FNN_{1}\) 的面积分别为 \(S_{1}\)，\(S_{2}\)，\(S_{3}\)，试判断 \(S_{2}^{2} = 4S_{1}S_{3}\) 是否成立，并证明你的结论.
\end{enumerate}

\bitmapfigure[max width=0.42\linewidth]{img/2009/hubei_liberal/q20_fig1.png}
\end{problem}

\begin{solution}
\begin{enumerate}
\item 令抛物线参数为 \(t\)，则其上的点可表示为 \(\left(\dfrac{p}{2}t^2,pt\right)\)，焦点为 \(F\left(\dfrac{p}{2},0\right)\)，准线为 \(x=-\dfrac{p}{2}\). 设 \(M\)，\(N\) 对应的参数分别为 \(t_1\)，\(t_2\)，则
\(M\left(\frac{p}{2}t_1^2,pt_1\right)\)，\(N\left(\frac{p}{2}t_2^2,pt_2\right)\)，\(M_1\left(-\frac{p}{2},pt_1\right)\)，\(N_1\left(-\frac{p}{2},pt_2\right)\).
过参数为 \(t_1\)，\(t_2\) 的两点的弦方程为 \((t_1+t_2)y=2x+pt_1t_2\). 因直线 \(MN\) 过焦点 \(F\)，代入 \(\left(\dfrac{p}{2},0\right)\) 得 \(t_1t_2=-1\). 于是
\(\overrightarrow{FM_1}=p(-1,t_1)\)，\(\overrightarrow{FN_1}=p(-1,t_2)\)，\(\overrightarrow{FM_1}\cdot\overrightarrow{FN_1}=p^2(1+t_1t_2)=0\).
所以 \(FM_1\perp FN_1\).
\item 由三角形面积公式，
\(S_1=\frac{p^2|t_1|(t_1^2+1)}{4}\)，\(S_3=\frac{p^2|t_2|(t_2^2+1)}{4}\)，\(S_2=\frac{p^2|t_1-t_2|}{2}\).
因为 \(t_1t_2=-1\)，所以 \(|t_1t_2|=1\)，并且
\[
(t_1^2+1)(t_2^2+1)=t_1^2t_2^2+t_1^2+t_2^2+1=t_1^2+t_2^2+2=(t_1-t_2)^2
\]
因此
\[
4S_1S_3=\frac{p^4}{4}(t_1^2+1)(t_2^2+1)=\frac{p^4(t_1-t_2)^2}{4}=S_2^2
\]
所以 \(S_2^2=4S_1S_3\) 成立.
\end{enumerate}
\end{solution}

\begin{problem}
已知关于 \(x\) 的函数 \(f(x) = -\frac{1}{3} x^{3} + bx^{2} + cx + bc\)，其导函数为 \(f'(x)\). 令 \(g(x) = |f'(x)|\)，记函数 \(g(x)\) 在区间 \([-1,1]\) 上的最大值为 \(M\).

\begin{enumerate}
\item 如果函数 \( f(x) \) 在 \( x = 1 \) 处有极值 \( -\frac{4}{3} \)，试确定 \( b\)，\( c \) 的值；
\item 若 \( |b| > 1 \)，证明对任意的 \( c \)，都有 \( M > 2 \)；
\item 若 \(M \geqslant k\) 对任意的 \(b\)，\(c\) 恒成立. 试求 \(k\) 的最大值.
\end{enumerate}
\end{problem}

\begin{solution}
\begin{enumerate}
\item 由 \(f\) 在 \(x=1\) 处有极值，得 \(f'(1)=0\)，即 \(-1+2b+c=0\)，所以 \(c=1-2b\). 又 \(f(1)=-\dfrac43\)，即
\[
-\frac13+b+c+bc=-\frac43
\]
代入 \(c=1-2b\)，得 \(\dfrac23-2b^2=-\dfrac43\)，所以 \(b^2=1\). 若 \(b=1\)，则 \(c=-1\)，且 \(f'(x)=-(x-1)^2\)，在 \(x=1\) 两侧函数均递减，不构成极值，舍去. 因而 \(b=-1\)，\(c=3\). 此时 \(f'(x)=-(x-1)(x+3)\)，在 \(x=1\) 左侧为正、右侧为负，确为极大值点.
\item 记 \(q(x)=f'(x)=-x^2+2bx+c\). 则 \(q(1)-q(-1)=4b\). 若区间 \([-1,1]\) 上始终有 \(|q(x)|\leq2\)，则
\[
|4b|=|q(1)-q(-1)|\leq|q(1)|+|q(-1)|\leq4
\]
这与 \(|b|>1\) 矛盾. 所以 \(q(1)\)，\(q(-1)\) 中至少有一个的绝对值大于 \(2\)，从而 \(M>2\).
\item 对任意的 \(b\)，\(c\)，有
\[
q(-1)+q(1)-2q(0)=-2
\]
因此 \(2\leq|q(-1)|+|q(1)|+2|q(0)|\leq4M\)，从而 \(M\geq\dfrac12\). 另一方面，取 \(b=0\)，\(c=\dfrac12\)，则 \(q(x)=\dfrac12-x^2\). 当 \(-1\leq x\leq1\) 时，\(-\dfrac12\leq q(x)\leq\dfrac12\)，故此时 \(M=\dfrac12\). 所以对任意 \(b\)，\(c\) 恒成立的不等式中，\(k\) 的最大值为 \(\dfrac12\).
\end{enumerate}
\end{solution}
